\documentclass[12pt]{article}

\input{../../_assets/preambulo.tex}


\begin{document}

    % 1. Foto de fondo
    % 2. Título
    % 3. Encabezado Izquierdo
    % 4. Color de fondo
    % 5. Coord x del titulo
    % 6. Coord y del titulo
    % 7. Fecha

    
    \input{../../_assets/portada}
    \portadaExamen{ffccA4.jpg}{Curvas y\\Superficies\\Examen II}{Curvas y Superficies. Examen II}{MidnightBlue}{-8}{28}{2026}{José Juan Urrutia Milán}

    \begin{description}
        \item[Asignatura] Curvas y Superficies.
        \item[Curso Académico] 2021/22.
        \item[Grado] Grado en Matemáticas.
        \item[Grupo] B.
        % \item[Profesor] ---.
        \item[Descripción] Convocatoria Ordinaria.
        \item[Fecha] 15 de junio de 2022.
        \item[Duración] 2 horas y media.
    \end{description}
    \newpage


    % ------------------------------------
    
    \begin{ejercicio}
        Realice los siguientes apartados:
        \begin{enumerate}[label=\alph*)]
            \item \textbf{[1 punto]} Sea $\alpha:\mathbb{R}\to\mathbb{R}^2$ la curva dada por:
                \begin{equation*}
                    \alpha(t) = ((1-2\sen(t))\cos(t), (1-\sen(t))\sen(t))
                \end{equation*}
                Prueba que $\alpha$ es una curva regular y calcula su curvatura. ¿Es $\alpha\big|_{[0,2\pi]}$ inyectiva?
            \item \textbf{[1 punto]} Sea $S$ una superficie orientable, $N$ una aplicación de Gauss en $S$ y $\alpha:I\to S$ una curva parametrizada por el arco. Consideramos $\hat{T}(s) = \alpha'(s)$, $\hat{N}(s) = N(\alpha(s))$ y $\hat{B}(s) = \hat{T}(s)\times \hat{N}(s)$. Se llama torsión geódesica de $\alpha$ a la función
                \begin{equation*}
                    \tau_g(s) = \langle \hat{B}'(s), \hat{N}(s) \rangle 
                \end{equation*}
                Demuestra que $\tau_g(s) = 0$ para todo $s\in I$ si y solo si $\alpha$ es una línea de curvatura de $S$, es decir $\alpha'(s)$ es una dirección principal de $S$ en $\alpha(s)$ para todo $s\in I$.
        \end{enumerate}
    \end{ejercicio}

    \begin{ejercicio}
        Sea $S = \{X(u,v) : u\in \left]0,\pi\right[, v\in \mathbb{R}\}$ donde
        \begin{equation*}
            X(u,v) = (v\cos(u), v\sen(u), u^2-2v+2)
        \end{equation*}
        y la curva $\alpha(t) = (\cos(t), \sen(t),t^2)$, $t\in \left]0,\pi\right[$. Se pide:
        \begin{enumerate}[label=\alph*)]
            \item \textbf{[0.5 puntos]} Demuestra que $X$ es una parametrización global de $S$.
            \item \textbf{[0.75 puntos]} Calcula la primera y segunda formas fundamentales de $S$.
            \item \textbf{[0.75 puntos]} Calcula la curvatura de Gauss y la curvatura media de $S$.
            \item \textbf{[0.75 puntos]} Prueba que el punto $p=(0,0,3)\in S$ y determinar el plano tangente de la superficie en ese punto. Calcula las curvaturas principales y las direcciones asintóticas en $p$, si las hubiera.
            \item \textbf{[0.75 puntos]} Comprueba que $\alpha$ es una curva asintótica de $S$.
            \item \textbf{[0.5 puntos]} Comprueba que $K(\alpha(t)) = -\tau_\alpha(t)^2$, donde $K(\alpha(t))$ denota la curvatura de Gauss de la superficie en el punto $\alpha(t)$ y $\tau_\alpha(t)$ es la torsión de la curva $\alpha$ en $t$.
        \end{enumerate}
    \end{ejercicio}

    \begin{ejercicio}
        Sea $X:\{(u,v)\in \mathbb{R}^2) : u^2+v^2 < 1\}\to \mathbb{R}^3$ una parametrización de una superficie $S$ tal que
        \begin{equation*}
            E(u,v) = G(u,v) = \frac{4}{{(1-u^2-v^2)}^{2}}, \qquad F(u,v) = 0
        \end{equation*}
        \begin{enumerate}[label=\alph*)]
            \item \textbf{[1 punto]} Calcula los símbolos de Christoffel de $S$ en esta parametrización.
            \item \textbf{[1 punto]} Calcula la curvatura de Gauss de $S$. ¿Existe una isometría local de $S$ en un plano?
        \end{enumerate}
    \end{ejercicio}

    \begin{ejercicio}
        Estudia de forma razonada las siguientes cuestiones:
        \begin{enumerate}[label=\alph*)]
            \item \textbf{[1 punto]} Sea $S$ una superficie y $p\in S$. Supongamos que cada entorno abierto de $p$ en $S$ contiene puntos a ambos lados del plano afín tangente a $S$ en $p.$ ¿Podemos afirmmar que $p$ es un punto hiperbólico?
            \item \textbf{[1 punto]} Una superficie $S$ se dice que es homogénea si para cada par de puntos $p_1$ y $p_2$ de $S$ existe una isometría $\Phi:S\to S$ tal que $\Phi(p_1) = p_2$. Prueba que toda superficie compacta, conexa y homogénea es una esfera.
        \end{enumerate}
    \end{ejercicio}

    \newpage
    \setcounter{ejercicio}{0}
    \noindent
    \textbf{Solución.}

    \begin{ejercicio}
        Realice los siguientes apartados:
        \begin{enumerate}[label=\alph*)]
            \item \textbf{[1 punto]} Sea $\alpha:\mathbb{R}\to\mathbb{R}^2$ la curva dada por:
                \begin{equation*}
                    \alpha(t) = ((1-2\sen(t))\cos(t), (1-\sen(t))\sen(t))
                \end{equation*}
                Prueba que $\alpha$ es una curva regular y calcula su curvatura. ¿Es $\alpha\big|_{[0,2\pi]}$ inyectiva?

                Calculamos $\alpha'$:
                \begin{align*}
                    \alpha'(t) &= (-2\cos^2(t) - \sen(t)(1-2\sen(t)), -\cos(t)\sen(t) + \cos(t)(1-\sen(t))) \\
                               &= (4\sen^2(t)-\sen(t)-2, \cos(t) -2\cos(t)\sen(t))
                \end{align*}
                Vemos que:
                \begin{equation*}
                    \alpha'(t) = 0 \quad\Longleftrightarrow\quad \left\{\begin{array}{l}
                        4\sen^2(t)-\sen(t)-2 = 0 \\
                        \cos(t)-2\cos(t)\sen(t) = 0
                    \end{array}\right.
                \end{equation*}
                Pero:
                \begin{equation*}
                    \cos(t)(1-2\sen(t)) = \cos(t)-2\cos(t)\sen(t) = 0 \quad\Longleftrightarrow\quad \left\{\begin{array}{l}
                        \cos(t) = 0 \\
                        ó \\
                        2\sen(t) = 1
                    \end{array}\right.
                \end{equation*}
                \begin{itemize}
                    \item Si $2\sen(t) = 1$ entonces:
                        \begin{equation*}
                            0 = 4\sen^2(t)-\sen(t)-2 = {(2\sen(t))}^{2}-\sen(t)-2
                        \end{equation*}
                        Multiplicando por 2:
                        \begin{equation*}
                            0 = 2{(2\sen(t))}^{2}-2\sen(t)-4 = 2-1-4 = -3
                        \end{equation*}
                        hemos llegado a una contradicción, este caso es imposible.
                    \item Si $\cos(t) = 0$ entonces $\sen(t) = \pm 1$.

                        Si $\sen(t) = 1$ entonces:
                        \begin{equation*}
                            0 = 4\sen^2(t) -\sen(t)-2 = 4-1-2 = 1
                        \end{equation*}
                        Imposible, y si $\sen(t) = -1$ entonces:
                        \begin{equation*}
                            0 = 4\sen^2(t) -\sen(t)-2 = 4+1-2 = 3
                        \end{equation*}
                        Imposible
                \end{itemize}
                Así $\alpha'(t)\neq 0 \quad \forall t\in \mathbb{R}$, por lo que $\alpha$ es regular.

                Como hemos visto que $\alpha$ es regular podemos aplicar la fórmula:
                \begin{equation*}
                    k_\alpha(t) = \frac{\det(\alpha'(t),\alpha''(t))}{|\alpha'(t)|^3}
                \end{equation*}
                Teníamos ya que $\alpha'(t) = (4\sen^2(t)-\sen(t)-2, \cos(t) -2\cos(t)\sen(t))$, de donde:
                \begin{equation*}
                    \alpha''(t) = (8\sen(t)\cos(t)-\cos(t),-\sen(t)+2\sen^2(t)-2\cos^2(t)) 
                \end{equation*}
                Basta sustituir en la fórmula de la curvatura.
            \item \textbf{[1 punto]} Sea $S$ una superficie orientable, $N$ una aplicación de Gauss en $S$ y $\alpha:I\to S$ una curva parametrizada por el arco. Consideramos $\hat{T}(s) = \alpha'(s)$, $\hat{N}(s) = N(\alpha(s))$ y $\hat{B}(s) = \hat{T}(s)\times \hat{N}(s)$. Se llama torsión geódesica de $\alpha$ a la función
                \begin{equation*}
                    \tau_g(s) = \langle \hat{B}'(s), \hat{N}(s) \rangle 
                \end{equation*}
                Demuestra que $\tau_g(s) = 0$ para todo $s\in I$ si y solo si $\alpha$ es una línea de curvatura de $S$, es decir $\alpha'(s)$ es una dirección principal de $S$ en $\alpha(s)$ para todo $s\in I$.\\

                En primer lugar, debemos darnos cuenta de que aunque $T(s) = \hat{T}(s)$, no tiene por qué ser $N(s) = \hat{N}(s)$, pues $\hat{N}$ depende de la orientación de $S$. No podemos usar las ecuaciones de Frenet en este ejercicio.

                Mediante doble implicación, comenzamos por la segunda, que es la que nos dislumbra cómo hacer el ejercicio:
                \begin{description}
                    \item [$\Longleftarrow )$] Supuesto que $\alpha'(s)$ es una dirección principal de $S$ en $\alpha(s)$ para todo $s\in I$, tenemos entonces que si $A_p$ es el endomorfismo de Weingarten de $S$ en $p$ entonces existe una aplicación $\lm(s):I\to\mathbb{R}$ de modo que:
                        \begin{equation*}
                            A_{\alpha(s)}\hat{T}(s) = \lm(s) \hat{T}(s)
                        \end{equation*}
                        Ahora, si tenemos buen ojo podemos observar que:
                        \begin{equation*}
                            A_{\alpha(s)}\hat{T}(s) = A_{\alpha(s)}\alpha'(s) = -(dN)_{\alpha(s)}(\alpha'(s)) = -(N\circ \alpha)'(s) = -\hat{N}'(s)
                        \end{equation*}
                        de donde $\hat{N}'(s) = -\lm(s)\hat{T}(s)$. Si ahora miramos la definición de $\tau_g$, tenemos $\hat{B}'$ multiplicado por $\hat{N}$, pero con la información superior solo sabemos que $\hat{N}'$ multiplicado por $\hat{B}$ es $0$. Para intercambiar las derivadas podemos usar la fórmula de la derivada de un producto:
                        \begin{equation*}
                            \tau_g(s) = \langle \hat{B}'(s), \hat{N}(s) \rangle = -\langle \hat{B}(s),\hat{N}'(s) \rangle + \cancel{(\langle \hat{B}(s),\hat{N}(s) \rangle )'} = 0
                        \end{equation*}
                        La cancelación se debea que $\hat{B}\perp \hat{N}$.
                    \item [$\Longrightarrow )$] Todos los pasos de la imiplicación anteriro son revertibles. Si $\tau_g(s) = 0$ tenemos entonces que:
                        \begin{equation*}
                            0 = \tau_g(s) = \langle \hat{B}'(s), \hat{N}(s) \rangle = -\langle \hat{B}(s),\hat{N}'(s) \rangle + \cancel{(\langle \hat{B}(s),\hat{N}(s) \rangle )'} 
                        \end{equation*}
                        Ahora, como $|N(\alpha(s))| = |\hat{N}(s)|^2 = 1$ por ser $N$ una aplicación de Gau\ss, derivando vemos que:
                        \begin{equation*}
                            0 = \langle \hat{N},\hat{N}' \rangle 
                        \end{equation*}
                        Por ser $\{\hat{T},\hat{N},\hat{B}\}$ una base de $\mathbb{R}^3$ deducimos que ha de existir $\lm:I\to\mathbb{R}$ de modo que:
                        \begin{equation*}
                            \hat{N}'(s) = \lm(s)\hat{T}(s)
                        \end{equation*}
                        Sin embargo, en la implicación anterior vimos que $\hat{N}'(s) = -A_{\alpha(s)}\hat{T}(s)$, por lo que acabamos de probar que $\alpha'(s) = \hat{T}(s)$ es una dirección principal de $S$ en $\alpha(s)$ para todo $s\in I$.
                \end{description}
        \end{enumerate}
    \end{ejercicio}

    \begin{ejercicio}
        Sea $S = \{X(u,v) : u\in \left]0,\pi\right[, v\in \mathbb{R}\}$ donde
        \begin{equation*}
            X(u,v) = (v\cos(u), v\sen(u), u^2-2v+2)
        \end{equation*}
        y la curva $\alpha(t) = (\cos(t), \sen(t),t^2)$, $t\in \left]0,\pi\right[$. Se pide:
        \begin{enumerate}[label=\alph*)]
            \item \textbf{[0.5 puntos]} Demuestra que $X$ es una parametrización global de $S$.

                Vemos en primer lugar que $X$ está definida en $U =\left]0,\pi\right[\times \mathbb{R}\subset\mathbb{R}^2$, que es un conjunto abierto. Vemos además que $X$ es una aplicación diferenciable. Falta ver que $X$ es inyectiva y que es un homeomorfismo sobre su imagen (es decir, que su aplicación inversa es continua).

                \begin{description}
                    \item [Inyectividad.] Si $(u_1,v_1),(u_2,v_2)\in U$ con $X(u_1,v_1) = X(u_2,v_2)$ tenemos entonces que:
                        \begin{equation*}
                            v_1\cos u_1 = v_2\cos u_2, \qquad v_1\sen u_1 = v_2\sen u_2
                        \end{equation*}
                        deducimos que $v_1^2 = v_2^2$. Distinguimos casos:
                        \begin{itemize}
                            \item Si $v_1=v_2$ como $\cos$ es biyectiva en $\left]0,\pi\right[$ obtenemos que $u_1 = u_2$.
                            \item Si $v_1=-v_2$ como $\sen u\neq 0$ para $u \in \left]0,\pi\right[$ esta situación es imposible\footnote{Pues tendríamos que $v_1\sen u_1 = v_2\sen u_2$.} a no ser que $v_1 = 0 = v_2$, pero ya queda tratada en el caso anterior.
                        \end{itemize}
                        Así, $X$ es inyectiva.
                    \item [Inversa continua.] Para $(x,y,z)\in X(U)$ buscamos $(x,y,z)\mapsto (u,v)$ que sea continua. Buscamos $v\in \mathbb{R}$ de manera que $v^2 = x^2+y^2$, por lo que las posibles elecciones son:
                        \begin{equation*}
                            \sqrt{x^2+y^2}, \qquad -\sqrt{x^2+y^2}
                        \end{equation*}
                        La observación clave es que $\sen(u)>0$ para $u\in \left]0,\pi\right[$, por lo que tenemos $sgn(v) = sgn(y)$. Así, tomamos:
                        \begin{equation*}
                            v = sgn(y)\sqrt{x^2+y^2}
                        \end{equation*}
                        para:
                        \begin{equation*}
                            sgn(y) = \left\{\begin{array}{ll}
                                -1 & \text{si\ } y<0 \\
                                0 & \text{si\ } y=0\\
                                1 & \text{si\ } y>0
                            \end{array}\right. 
                        \end{equation*}
                        Y observamos que $v = sgn(y)\sqrt{x^2+y^2}$ es una función continua, porque de $y\to 0^+$ por ser $\sen(u)>0$ debemos tener $v\to 0^+$, de donde $x\to 0^+$ ya la función es continua para $y = 0$. Una vez tenemos $v$ podemos obtener $u$ despejando en la tercera ecuación ($u\in \left]0,\pi\right[$):
                        \begin{equation*}
                            z = u^2-2v+2 \quad\Longrightarrow\quad 0 < u^2 = z+2v+2
                        \end{equation*}
                        de donde:
                        \begin{equation*}
                            u = \sqrt{z+2v+2}
                        \end{equation*}
                        donde $v$ era una aplicación continua, por lo que $u$ es también continua.
                \end{description}
                Finalmente, falta ver también que $(dX)_q$ es inyectiva para todo $q\in U$.
            \item \textbf{[0.75 puntos]} Calcula la primera y segunda formas fundamentales de $S$.

                Lo haremos a partir de la parametrización $X$, por lo que debemos calcular $E,G,F,e,g,f$. Para ello:
                \begin{gather*}
                    X_u = (-v\sen(u), v\cos(u), 2u), \qquad X_v = (\cos(u),\sen(u),-2) \\
                    X_{uu} = (-v\cos(u), -v\sen(u), 2), \quad X_{uv} = (-\sen(u),\cos(u),0), \quad X_{vv} = (0,0,0)
                \end{gather*}
                Así:
                \begin{align*}
                    E &= |X_u|^2 = v^2 + 4u^2, \qquad G = |X_v|^2 = 1+4 = 5 \\
                    F &= \langle X_u,X_v \rangle = -4u
                \end{align*}
                de donde:
                \begin{equation*}
                    EG -F^2 = 5v^2 + 20u^2 -16u^2 = 5v^2 + 4u^2
                \end{equation*}
                seguimos calculando:
                \begin{align*}
                    e &= \frac{\left|\begin{array}{ccc}
                        -v\sen(u) & v\cos(u) & 2u \\
                        \cos(u) & \sen(u) & -2 \\
                        -v\cos(u) & -v\sen(u) & 2
                    \end{array}\right|}{\sqrt{EG-F^2}}
                \end{align*}
            \item \textbf{[0.75 puntos]} Calcula la curvatura de Gauss y la curvatura media de $S$.
            \item \textbf{[0.75 puntos]} Prueba que el punto $p=(0,0,3)\in S$ y determinar el plano tangente de la superficie en ese punto. Calcula las curvaturas principales y las direcciones asintóticas en $p$, si las hubiera.
            \item \textbf{[0.75 puntos]} Comprueba que $\alpha$ es una curva asintótica de $S$.
            \item \textbf{[0.5 puntos]} Comprueba que $K(\alpha(t)) = -\tau_\alpha(t)^2$, donde $K(\alpha(t))$ denota la curvatura de Gauss de la superficie en el punto $\alpha(t)$ y $\tau_\alpha(t)$ es la torsión de la curva $\alpha$ en $t$.
        \end{enumerate}
    \end{ejercicio}

    \begin{ejercicio}
        Sea $X:\{(u,v)\in \mathbb{R}^2) : u^2+v^2 < 1\}\to \mathbb{R}^3$ una parametrización de una superficie $S$ tal que
        \begin{equation*}
            E(u,v) = G(u,v) = \frac{4}{{(1-u^2-v^2)}^{2}}, \qquad F(u,v) = 0
        \end{equation*}
        \begin{enumerate}[label=\alph*)]
            \item \textbf{[1 punto]} Calcula los símbolos de Christoffel de $S$ en esta parametrización.
            \item \textbf{[1 punto]} Calcula la curvatura de Gauss de $S$. ¿Existe una isometría local de $S$ en un plano?
        \end{enumerate}
    \end{ejercicio}

    \begin{ejercicio}
        Estudia de forma razonada las siguientes cuestiones:
        \begin{enumerate}[label=\alph*)]
            \item \textbf{[1 punto]} Sea $S$ una superficie y $p\in S$. Supongamos que cada entorno abierto de $p$ en $S$ contiene puntos a ambos lados del plano afín tangente a $S$ en $p.$ ¿Podemos afirmmar que $p$ es un punto hiperbólico?
            \item \textbf{[1 punto]} Una superficie $S$ se dice que es homogénea si para cada par de puntos $p_1$ y $p_2$ de $S$ existe una isometría $\Phi:S\to S$ tal que $\Phi(p_1) = p_2$. Prueba que toda superficie compacta, conexa y homogénea es una esfera.
        \end{enumerate}
    \end{ejercicio}


\end{document}
